\documentclass[../DoAn.tex]{subfiles}
\begin{document}

\section{Khảo sát hiện trạng và đánh giá tương tự}
\label{section:2.1}
Trên thị trường hiện nay có nhiều tựa game hành động góc nhìn thứ ba (TPS) và thứ nhất (FPS) nổi bật kết hợp cơ chế thao tác không gian, điển hình như dòng game giải đố \textit{Portal} hay tựa game bắn súng nhiều người chơi \textit{Splitgate}. Tuy nhiên, phần lớn đều tập trung cho giải đố hoặc nhiều người chơi online, số lượng các trò chơi bắn súng góc nhìn thứ ba (TPS) kết hợp cổng không gian theo hướng chơi đơn (PvE) với cốt truyện khoa học viễn tưởng kịch tính vẫn còn nhiều khoảng trống để khai thác. Cốt truyện của dự án xoay quanh cuộc chiến chống lại những kẻ thù lội ngược dòng thời gian từ tương lai nhằm mục đích thay đổi quá khứ (chi tiết cốt truyện sẽ được phân tích ở Chương Thiết kế trò chơi). Nhờ đó, nhân vật chính sử dụng linh hoạt khả năng mở cổng dịch chuyển (Portal) ở góc nhìn thứ ba để tạo lợi thế di chuyển, khai thác các góc bắn bất ngờ nhằm hạ gục kẻ địch.

Về mặt nguồn lực và chi phí mức hiện trạng, dự án định hướng tận dụng tối đa nền tảng đồ họa sức mạnh cao Unreal Engine 5.6. Thay vì thiết kế nội dung 3D từ con số không, dự án sử dụng tài nguyên (asset) theo định dạng miễn phí hoặc chi phí thấp trên cửa sổ tài nguyên Fab/Epic Games. Bằng cách tự tối ưu thông số, tùy chỉnh các hiệu ứng hình ảnh (VFX), can thiệp hệ thống ánh sáng, đồ họa của trò chơi vẫn được duy trì ở mức tiệm cận của một sản phẩm chất lượng cao sắc nét, nhưng tổng mức đầu tư tài chính cho cấu thành tài nguyên mới chỉ dưới 1.500.000 VNĐ. Điểm mạnh này đánh bại nhiều dự án có ngân sách tương tự, chứng minh tính khả thi của việc tạo ra tựa game tốt song song với việc tối ưu chi phí thấp.

\section{Mục đích của trò chơi}
\label{section:2.2}
Mục tiêu cốt lõi của dự án là cung cấp một sản phẩm máy tính có tính tương tác cao dành cho người dùng trải nghiệm không gian giải trí cao cấp. Bên cạnh việc thỏa mãn thư giãn đơn thuần, trò chơi hướng tới việc là một môi trường rèn luyện phản xạ (reflexes) ở nhịp độ cao. Thêm vào đó, thông qua cơ chế cổng không gian đa hướng mới mẻ, trò chơi buộc game thủ phải định hình tư duy logic không gian tốt, liên tục phân tích địa hình nhằm thiết lập các kế hoạch lên góc bắn (tactical planning) chủ động.

\section{Định hướng sử dụng}
\label{section:2.3}
Dự án được ứng dụng dưới dạng một trò chơi độc lập ngoại tuyến (Offline Stand-alone Game) dành cho hệ máy tính (PC). Hướng sử dụng chính đặt tại các thời điểm và nhu cầu thao tác kỹ năng cá nhân của người chơi. Mặc dù nhiều trò chơi mang lại hiệu năng giải tỏa căng thẳng (stress) tâm lý, nhưng do vấn đề tâm lý không có cơ sở khoa học y tế đánh giá hiệu chuẩn chính xác trong quy mô hệ kỹ thuật, nên đồ án loại bỏ định hướng định lượng mơ hồ này. Định hướng thực tế của đồ án là đóng vai trò như một sản phẩm thương mại giải trí tiêu chuẩn trong thị trường game thủ, đồng thời là một dự án nghiên cứu khả năng tinh chỉnh đồ họa UE 5.6 đối với mô hình sản xuất vi mô ngân sách dưới 1.5 triệu đồng (Low-budget Indie).

\section{Cơ sở của định hướng sử dụng}
\label{section:2.4}
Tính phù hợp vững chắc trong định hướng phát triển của trò chơi được phân tích trên tiến trình thực tiễn của công nghệ và thể loại trò chơi đã được thị trường công nhận. Việc kết nối sự tự do của cơ chế thiết lập \textit{Portal} và tiết tấu bắn súng góc nhìn thứ ba (TPS) luôn đem lại sự cuốn hút lôi kéo người chơi rất lớn - biểu hiện thực tế từ lượt theo dõi bền vững trên các thị trường thương mại và cộng đồng.

Bên cạnh đó, chiến lược áp dụng các Asset có sẵn để cấu tạo trên nền Unreal Engine công nghệ mới mà không cần đầu tư sản xuất mô hình khổng lồ đang là quy chuẩn của các nhà phát triển độc lập (Indies) chuyên nghiệp toàn cầu nhằm đi đường vòng khéo léo qua bài toán chi phí đắt đỏ của quá khứ. Đây là một cơ sở thực chứng và khoa học củng cố cho sự chắc chắn và toàn vẹn của hướng phát triển này để mang đến đồ án có chất lượng kỹ thuật cực kỳ sâu không kém các đội nhóm lớn.

%%%%%%%%%%%%%%%%%%%%%%%%%%%%%%%%%%%

\end{document}